\subsection{Угол между вектором и плоскостью}

Найдите угол между вектором~$v=\ee(1, -5, -3, -5)^\top$ и плоскостью~$\pi$, заданной системой уравнений
\begin{equation}
	\label{eq:num:b:plane_system}
	\begin{system}
		-4x + 3y + 2z + t 
		& = 0,
		\\
		x - 2y + z
		& = 0.
	\end{system}
\end{equation}

\begin{proof}
	В нашем случае плоскость задана как ядро матрицы
	\[
		B :=
		\begin{pmatrix}
			-4 & 3 & 2 & 1 \\
			1 & -2 & 1 & 0
		\end{pmatrix}
		.
	\]
	Чтобы это увидеть, заметим, что система~\eqref{eq:num:b:plane_system} эквивалентна уравнению
	\[
		Bw = 0
		,
	\]
	где $w = (x, y, z, t)^\top$.

	Если бы плоскость~$\pi$ была задана как двумерная плоскость в трёхмерном пространстве, то найти косинус угла~$\theta$ между~$v$ и~$\pi$ можно было бы, найдя угол~$\pi/2 - \theta$ через скалярное произедение~$\spr{v}{n}$, где~$n$ --- нормаль к~$\pi$. Тогда
	\[
		\cos
		\left( 
			\frac{\pi}{2} -
			\theta
		\right)
		=
		\sin{\theta}
		=
		\frac{\spr{v}{n}}{\norm{v} \cdot \norm{n}}
		=
		\spr{e_v}{e_n}
		,
	\]
	где~$e_v := v/\norm{v}$, $e_n := n/\norm{n}$. Получилось, что
	\begin{equation}
		\label{eq:num:b:sin:proj}
		\sin{\theta} = \pr_{e_n}{e_v} = \pr_n{e_v}
		,
	\end{equation}
	Мы воспользовались бы тем, что ортогональное дополнение до~$\pi$ одномерно.

	Но в нашем случае ортогональное дополнение до плоскости~$\pi$ двумерно, как и сама плоскость~$\pi$, и поэтому надо искать угол между одномерным подпространством~$\anspan{\vec{v}}$ и двумерным подпространством~$\pi$ либо его ортогональным дополнением~$\pi^{\bot}$.
	
	Однако в равенстве~\eqref{eq:num:b:sin:proj} видно, что синус угла задаётся проекцией на ортогональное дополнение. Эту идею можно использовать и в многомерии.

	Можно найти и проекцию на плоскость непосредственно, но для этого придётся найти базис~$\pi$.

	Но в виде~\eqref{eq:num:b:plane_system} нам \emph{уже дан} базис в ортогональном дополнении~$\pi^\bot$. 

	Координаты базисных векторов из~$\pi^\bot$ --- это строки матрицы~$B$, так как уравнение~$Bw = 0$ задаёт такие векторы~$w$, для которых скалярное произведение $i$-ой строки~$b^i$ матрицы~$B$ на~$w$ равно~$0$, то есть задаёт все векторы~$w$, которые ортогональны всем строкам~$B$. Но векторы~$w$, удовлетворяющие уравнению~$Bw = 0$, --- это в точности векторы плоскости~$\pi$. Заметим, что~$\rank{B} = 2$ --- максимально возможный. Поэтому строки~$B$ линейно независимы. Поэтому они образуют координаты базисных векторов ортогонального дополнения к~$\pi$.

	\begin{wrapfigure}[10]{R}{0.35\linewidth}
		\vspace{-1.2\baselineskip}
		\centering
		\includegraphics{pics/triangle.pdf}
		\caption{Углы между проекциями вектора~$v$.}
		\label{pic:triangle}
	\end{wrapfigure}

	Итак, найдём проекцию~$v$ на~$\pi^\bot$ (обозначим её $v_\bot$), а затем из~$v$ и~$v_\bot$ найдём длину проекции~$v_\pi$ вектора~$v$ на~$\pi$. Для этого воспользуемся многомерной теоремой Пифагора. Так как~$v_\pi \in \pi$ и $v_\bot \in \pi^\bot$, а $\pi \perp \pi^\bot$, то~$v_\pi \perp v_\bot$, и поэтому
	\[
		\norm{v}^2 = \norm{v_\pi}^2 + \norm{v_\bot}^2.
	\]

	Нарисовав прямоугольный треугольник (см.~\picref{pic:triangle}) с катетами~$\norm{v_\pi}$ и~$\norm{v_\bot}$ и гипотенузой~$\norm{v}$, а затем обозначив буквой~$\theta$ угол между~$v$ и~$v_\pi$, а буквой~$\alpha$ --- угол между~$v$ и~$v_\bot$, увидим, что
	\[
		\cos{\alpha} = \sin{\theta} =
		\frac{\norm{v_\bot}}{\norm{v}}
		.
	\]
	Отсюда
	\begin{equation}
		\label{eq:num:b:alpha}
		\alpha =
		\arccos
		\left( 
			\frac{\norm{v_\bot}}{\norm{v}}
		\right)
		.
	\end{equation}
	До этой формулы и нужно дойти.


	Пусть~$\nn = (n_1, n_2)$ --- базис в~$\pi^\bot$:
	\[
		\nn = \ee N = \ee B^\top =
		\ee
		\begin{pmatrix}
			-4 & 1 \\
			3 & -2 \\ 
			2 & 1 \\
			1 & 0
		\end{pmatrix}
		.
	\]
	\begingroup
		\color{gray}
		Теперь найдём
		\[
			v_\bot =
			\pr_{n_1}{v} +
			\pr_{n_2}{v}
			=
			\vproj{v}{n_1} +
			\vproj{v}{n_2}
			.
		\]
		(Эти формулы не сработают, так как~$\nn$ --- не ортонормированный базис.)
		
		Можно либо привести базис~$\nn$ к ортонормированному ортогонализацией Грама--Шмидта, либо использовать рассуждения, изложенные ниже.
	\endgroup

	Пусть
	\[
		v = v_\top + v_\pi
	\]
	и $v_\pi \bot n_i$ для~$i \in \{1,2\}$.
	Мы хотим найти разложение
	\[
		v_\bot = a^1 n_1 + a^2 n_2
		.
	\]
	Для этого умножим обе стороны этого равенства скалярно на~$n_i$:
	\begin{equation}
		\label{eq:num:b:v_bot_system}
		\begin{system}
			\spr{v_\bot}{n_1} &=
			a^1 \spr{n_1}{n_1} +
			a^2 \spr{n_2}{n_1},
			\\
			\spr{v_\bot}{n_2} &=
			a^1 \spr{n_1}{n_2} +
			a^2 \spr{n_2}{n_2}.
		\end{system}
	\end{equation}
	Далее заметим, что
	\[
		\spr{v}{n_i} =
		\spr{v_\bot + v_\pi}{n_i} =
		\spr{v_\bot}{n_i} +
		\spr{v_\pi}{n_i} =
		\spr{v_\bot}{n_i}
		.
	\]
	Поэтому~\eqref{eq:num:b:v_bot_system} можно переписать в иде
	\begin{equation}
		\label{eq:num:b:v_system}
		\begin{system}
			\spr{v}{n_1} &=
			a^1 \spr{n_1}{n_1} +
			a^2 \spr{n_2}{n_1},
			\\
			\spr{v}{n_2} &=
			a^1 \spr{n_1}{n_2} +
			a^2 \spr{n_2}{n_2},
		\end{system}
	\end{equation}
	причём скалярное произведение всюду берётся евклидовым (с единичной матрицей Грама). Перепишем~\eqref{eq:num:b:v_system} в матричном виде:
	\begin{equation}
		\label{eq:num:b:v_matrix}
		\Vect{
			\spr{v}{n_1}
			\\
			\spr{v}{n_2}
		}
		=
		\begin{pmatrix}
			\spr{n_1}{n_1} & \spr{n_2}{n_1} \\
			\spr{n_1}{n_2} & \spr{n_2}{n_2}
		\end{pmatrix}
		\Vect{ a^1 \\ a^2 }
		.
	\end{equation}
	Это система линейных уравнений на~$a^i$. Скалярные произведения нам уже известны (их нужно только посчитать). Заметим, что матрица скалярных произведений --- это матрица~$N^\top N = B B^\top$. Посчитаем:
	\begin{align*}
		N^\top N =
		B B^\top &=
		\begin{pmatrix}
			-4 & 3 & 2 & 1 \\
			1 & -2 & 1 & 0
		\end{pmatrix}
		\begin{pmatrix}
			-4 & 1 \\
			3 & -2 \\
			2 & 1 \\
			1 & 0
		\end{pmatrix}
		% \nlinea{=}
		=
		\begin{pmatrix}
			30 & -8 \\
			-8 & 6
		\end{pmatrix}
		=
		\begin{pmatrix}
			\spr{n_1}{n_1} & \spr{n_2}{n_1} \\
			\spr{n_1}{n_2} & \spr{n_2}{n_2}
		\end{pmatrix}
		.
	\end{align*}

	Осталось найти скалярные произведения~$\spr{v}{n_1}$ и~$\spr{v}{n_2}$:
	\begin{align*}
		\spr{v}{n_1} &=
		\spr{
			\ee
			\Vect{1 \\ -5 \\ -3 \\ -5}
		}{
			\ee
			\Vect{-4 \\ 3 \\ 2 \\ 1}
		}
		=
		-4 - 15 - 6 - 5 = -30;
		\\
		\spr{v}{n_1} &=
		\spr{
			\ee
			\Vect{1 \\ -5 \\ -3 \\ -5}
		}{
			\ee
			\Vect{1 \\ -2 \\ 1 \\ 0}
		}
		=
		8
		.
	\end{align*}
	Итак, уравнение~\eqref{eq:num:b:v_matrix} принимает вид
	\begin{align*}
		\Vect{-30 \\ 8}
		&=
		\begin{pmatrix}
			30 & -8 \\
			-8 & 6
		\end{pmatrix}
		\Vect{a^1 \\ a^2}
		\\
		\Vect{-15 \\ 4}
		&=
		\underbrace{
			\begin{pmatrix}
				15 & -4 \\
				-4 & 3
			\end{pmatrix}
		}_{G}
		\Vect{a^1 \\ a^2}
	\end{align*}
	(в левой и правой части можно вынести множитель~$1/2$). Найдём
	\[
		G^{-1} =
		\frac{1}{\det{G}}\,G^\vee,
	\]
	где~$G^\vee$ --- матрица алгебраических дополнений. Считаем определитель
	\[
		\det{G} =
		15 \cdot 3 - 16 =
		(16-1) \cdot 3 - 16
		=
		16 \cdot 3 - 3 - 16 =
		16 \cdot 2 - 3 =
		32 - 3 = 29
	\]
	и матрицу алгебраических дополнений:
	\[
		G^\vee
		=
		\begin{pmatrix}
			\det(3) & -\det(-4) \\
			-\det(-4) & \det(15)
		\end{pmatrix}
		=
		\begin{pmatrix}
			3 & 4 \\
			4 & 15
		\end{pmatrix}
		.
	\]
	Тогда
	\[
		G^{-1}
		=
		\frac{1}{29}
		\begin{pmatrix}
			3 & 4 \\
			4 & 15
		\end{pmatrix}
		.
	\]
	Поэтому
	\[
		\Vect{a^1 \\ a^2} =
		\frac{1}{29}
		\begin{pmatrix}
			3 & 4 \\
			4 & 15
		\end{pmatrix}
		\Vect{-15 \\ 4}
		=
		\frac{1}{29}
		\Vect{-29 \\ 0}
		=
		\Vect{-1 \\ 0}
		.
	\]
	Итак,
	\[
		v_\bot =
		\Vect{
			n_1 & n_2
		}
		\Vect{
			a^1 \\ a^2
		}
		=
		\Vect{
			n_1 & n_2
		}
		\Vect{ -1 \\ 0 }
		=
		-n_1
		=
		\ee
		\Vect{4 \\ -3 \\ -2 \\ -1}
		.
	\]

	Далее найдём длины:
	\begin{align*}
		\norm{v}^2 &=
		1 + (-5)^2 + (-3)^2 + (-5)^2 =
		1 + 25 + 9 + 25 = 60;
		\\
		\norm{v_\bot}^2 &=
		4^2 + (-3)^2 + (-2)^2 + (-1)^2 =
		16 + 9 + 4 + 1 =
		30.
	\end{align*}
	Тогда по формуле~\eqref{eq:num:b:alpha}
	\[
		\alpha =
		\arccos
		\left( 
			\frac{\norm{v_\bot}}{\norm{v}}
		\right)
		=
		\arccos
		\left( 
			\frac{\sqrt{30}}{\sqrt{60}}
		\right)
		=
		\arccos
		\left( 
			\frac{1}{\sqrt{2}}
		\right)
		=
		\frac{\pi}{4}
		.
	\]
	Но~$\alpha$ --- угол между~$v$ и~$v_\bot$; тогда искомый угол~$\theta = \pi/2 - \alpha = \pi/4$.
\end{proof}

